% !PRINTED 45
\weberpage{45}

Es werden dann auch die durch (4), (5) bestimmten Werte $x$
und $\sqrt{f(x)}$ Funktionen dieser Veränderlichen, und (5) läßt sich
differentiieren.

Bezeichnen wir mit $\delta$ die Differentiation nach den in den
Koeffizienten von $P$, $Q$ vorkommenden Veränderlichen, wobei $x$
als konstant gilt, mit $dx$ das entsprechende Differential von $x$,
wenn $x$ durch (4) oder (5) als Funktion dieser Koeffizienten
bestimmt ist, so ergibt die Differentiation von (4)
\[ 2P\delta P - 2Q\delta Qf(x) + F'(x)\,dx = 0, \]
woraus mit Benutzung von (5):
\weq{6}{\dfrac{Q\,\delta P - P\,\delta Q}{F'(x)} = \dfrac{dx}{2\sqrt{f(x)}},}
und da der Zähler auf der linken Seite vom Grade $2m-2$,
$F(x)$ vom Grade $2m$ ist, so läßt sich die Formel (2) anwenden,
und es folgt
\weq{7}{\sum \dfrac{dx}{\sqrt{f(x)}} = 0,}
wenn die Summe auf sämtliche Wurzeln der Gleichung (5) erstreckt
 ist.

Wir wenden dieses Theorem auf den einfachsten Fall, nämlich
 $m=2$, an und erhalten dann folgenden Satz:

\gesperrt{Wenn}
\weq{8}{x_1, \sqrt{f(x_1)}; \quad x_2, \sqrt{f(x_2)}; \quad x_3, \sqrt{f(x_3)}; \quad x_4, \sqrt{f(x_4)}}
\gesperrt{vier Wertepaare sind, für die irgend eine Funktion der
Form}
\weq{9}{a + bx + cx^2 + \sqrt{f(x)}}
\gesperrt{verschwindet, so sind diese Wertepaare nicht gänzlich
voneinander unabhängig, sondern es besteht zwischen
ihnen die Differentialgleichung}
\weq{10}{\dfrac{dx_1}{\sqrt{f(x_1)}} + \dfrac{dx_2}{\sqrt{f(x_2)}} + \dfrac{dx_3}{\sqrt{f(x_3)}} + \dfrac{dx_4}{\sqrt{f(x_4)}} = 0.}

Die Abhängigkeit dieser vier Wertepaare, also eine Integration
 der Differentialgleichung (10), kann aber auch in algebraischer
 Weise ausgedrückt werden, indem man die Funktion
(9) für $x = x_1, x_2, x_3, x_4$ gleich Null setzt und dann $a, b, c$
eliminiert. Man erhält so die Determinantengleichung:
%CONT>>
